\chapter{Utilisation de ROS en réseau}

Il est fortement recommandé d'utiliser un réseau privé, sur lequel il est possible d'attribuer des adresses IP fixes (routeurs, réseau INSA-Iot). S'assurer que chaque apparail du réseau ait le même \textit{ROS\_DOMAIN\_ID}. Pour cela, utiliser les commandes suivantes sur chaque appareil : 
\begin{lstlisting}[style=commandstyle]
echo export ROS_DOMAIN_ID=0 >> ~/.bashrc
source ~/.bashrc
\end{lstlisting}

\noindent Connecter ensuite tous les appareils au même réseau, tous les topics seront alors accessibles à tous les appareils du réseau. Cela permet de faire un contrôle centralisé des différents robots. Par exemple, pour un essaim, un ordinateur central pourra directement publier sur les topics \textit{cmd\_vel} attribués à chaque module.

